\section*{Capítulo \ref{ch:b1:photosynthesis} --- Fotosíntesis y autotrofía}

\begin{solution}{exo:b1:photosynthesis:1}
Las membranas externa e interna de la envoltura, y la membrana tilacoidal
del interior. \omterm{prop:b1:photosynthesis:lightreactions}{Reacciones luminosas}: en la membrana tilacoidal (con protones
bombeados a la luz tilacoidal). \omterm{def:b1:photosynthesis:fixation}{Ciclo de Calvin}: en el \omterm{def:b1:photosynthesis:chloroplast}{estroma}.
\end{solution}

\begin{solution}{exo:b1:photosynthesis:2}
$2\,\mathrm{H_2O} + 2\,\mathrm{NADP^+} + 3\,\mathrm{ADP} + 3\,\mathrm{P_i}
\to \mathrm{O_2} + 2\,\mathrm{NADPH} + 2\,\mathrm{H^+} + 3\,\mathrm{ATP}$
con ocho fotones. El oxígeno procede del agua: las algas a las que se da
$\mathrm{H_2^{18}O}$ liberan $\mathrm{^{18}O_2}$; si se les da
$\mathrm{C^{18}O_2}$, no (Ruben y Kamen).
\end{solution}

\begin{solution}{exo:b1:photosynthesis:3}
\omterm{def:b1:photosynthesis:fixation}{Fijación} (\omterm{def:b1:photosynthesis:fixation}{RuBisCO}), reducción (a G3P) y regeneración (de la RuBP). Por cada
$\mathrm{CO_2}$: 3 \omterm{def:b1:water-small-molecules:atp}{ATP} y 2 NADPH.
\end{solution}

\begin{solution}{exo:b1:photosynthesis:4}
Cerca de \qty{430}{nm} (azul) y de \qty{660}{nm} (rojo). En torno a
\qty{500}{nm} los \omterm{def:b1:photosynthesis:pigments}{carotenoides} absorben y transmiten la energía a la
\omterm{def:b1:photosynthesis:pigments}{clorofila}, de modo que se libera oxígeno aunque la \omterm{def:b1:photosynthesis:pigments}{clorofila} apenas absorba
allí.
\end{solution}

\begin{solution}{exo:b1:photosynthesis:5}
La luz roja lejana (\qty{700}{nm}) por sí sola impulsa poca fotosíntesis;
añadida a la luz de \qty{680}{nm} da más que la suma de las dos velocidades.
Dos sistemas de pigmentos con máximos de absorción distintos deben cooperar
en serie, y cada uno necesita sus propios fotones: el \omterm{def:b1:photosynthesis:zscheme}{fotosistema} I (P700) y
el \omterm{def:b1:photosynthesis:zscheme}{fotosistema} II (P680).
\end{solution}

\begin{solution}{exo:b1:photosynthesis:6}
$0.18\times 96\,500 = \qty{17.4}{kJ/mol}$ de protones; $50/17.4 = 2.9$
protones como mínimo; la sintasa usa cuatro, es decir, un rendimiento del
\qty{72}{\%}.
\end{solution}

\begin{solution}{exo:b1:photosynthesis:7}
Demuestra que un gradiente de protones a través de la membrana tilacoidal
basta para impulsar la síntesis de \omterm{def:b1:water-small-molecules:atp}{ATP}, sin luz y sin transporte de
electrones: el acoplamiento es a través del gradiente. No demuestra por sí
solo que las \omterm{prop:b1:photosynthesis:lightreactions}{reacciones luminosas} fabriquen \omterm{def:b1:water-small-molecules:atp}{ATP} así en la \omterm{def:b1:flowering-plant-organization:organs}{hoja} (para eso
hacen falta el desacoplante y las medidas de pH). Con un desacoplante el
gradiente se derrumba antes de que la sintasa pueda usarlo: no hay \omterm{def:b1:water-small-molecules:atp}{ATP}.
\end{solution}

\begin{solution}{exo:b1:photosynthesis:8}
$3\times 5 + 3\times 1 = 18$ carbonos; seis 3-PGA $= 18$; seis G3P $= 18$;
sale un G3P (3) y cinco G3P (15) regeneran tres RuBP (15): $3 +
15 = 18$.
Cuadra.
\end{solution}

\begin{solution}{exo:b1:photosynthesis:9}
Sin $\mathrm{CO_2}$ la \omterm{def:b1:photosynthesis:fixation}{RuBisCO} se detiene: la RuBP deja de consumirse (sube)
y el 3-PGA deja de fabricarse (baja) mientras la luz lo sigue reduciendo.
Sin luz no hay \omterm{def:b1:water-small-molecules:atp}{ATP} ni NADPH: el 3-PGA deja de reducirse (sube) y la RuBP
deja de regenerarse (baja) mientras la fijación continúa brevemente.
\end{solution}

\begin{solution}{exo:b1:photosynthesis:10}
$E = hc/\lambda = 6.63\times 10^{-34}\times 3\times 10^8/6.8\times
10^{-7} = \qty{2.92e-19}{J}$; por mol, \qty{176}{kJ}. 48 fotones:
\qty{8450}{kJ}; rendimiento $2870/8450 = \qty{34}{\%}$. Un cultivo pierde la
luz que no absorbe (media parte del espectro, reflexión, transmisión, huecos
entre plantas), la fotorrespiración, la respiración de \omterm{def:b1:flowering-plant-organization:organs}{hojas}, \omterm{def:b1:flowering-plant-organization:organs}{tallos} y
raíces, la saturación lumínica de las \omterm{def:b1:enzymes:enzyme}{enzimas} a pleno sol, y las estaciones
sin \omterm{def:b1:flowering-plant-organization:organs}{hojas}.
\end{solution}

\begin{solution}{exo:b1:photosynthesis:11}
Por cada 100 carboxilaciones, 33 oxigenaciones que liberan $16.7$ carbonos:
neto $83.3$, una pérdida del \qty{17}{\%}. Coste de evitarlo: 2 \omterm{def:b1:water-small-molecules:atp}{ATP}
$\approx
0.2$ carbonos por $\mathrm{CO_2}$, es decir, un \qty{20}{\%}; el
diseño C4 sale a cuenta solo cuando la pérdida supera la quinta parte, en
condiciones cálidas y secas, no en las frescas.
\end{solution}

\begin{solution}{exo:b1:photosynthesis:12}
Ambas usan membranas que bombean protones con la energía del transporte de
electrones y una \omterm{thm:b1:photosynthesis:chemiosmosis}{ATP sintasa} que gasta el gradiente: el mecanismo es el
mismo, y las sintasas son homólogas. Pero los electrones corren en sentidos
opuestos --- del agua al NADPH en el \omterm{def:b1:photosynthesis:chloroplast}{cloroplasto}, elevados por la luz; del
NADH al oxígeno en la \omterm{def:b1:eukaryotic-cell:mitochondrion}{mitocondria}, cayendo
(\cref{ch:b1:respiration-fermentation}) --- y los ciclos del carbono no son
inversos el uno del otro: el \omterm{def:b1:photosynthesis:fixation}{ciclo de Calvin} y el de Krebs no comparten
ninguna \omterm{def:b1:enzymes:enzyme}{enzima}, y la reducción del $\mathrm{CO_2}$ usa NADPH mientras que su
liberación usa $\mathrm{NAD^+}$. Las ecuaciones globales sí son inversas; la
maquinaria es una herencia común usada en dos sentidos, con química distinta
en el extremo del carbono.
\end{solution}

\begin{solution}{pb:b1:photosynthesis:1}
\textbf{1.} $hc/\lambda = \qty{2.92e-19}{J}$; $\times N_A = \qty{176}{kJ/mol}$.
\textbf{2.} \qty{700}{nm}: \qty{171}{kJ/mol}; \qty{450}{nm}:
\qty{266}{kJ/mol}. Un fotón azul excita la \omterm{def:b1:photosynthesis:pigments}{clorofila} a un estado superior
que decae en picosegundos al mismo estado excitado más bajo al que llega un
fotón rojo; el exceso se pierde como calor, y la química ve un solo fotón en
cualquiera de los dos casos.
\textbf{3.} $1/2.92\times 10^{-19} = \num{3.4e18}$ fotones por julio.
\textbf{4.} $\Delta E = 0.82 - (-0.32) = \qty{1.14}{V}$: $\Delta G =
96\,500\times 1.14 = \qty{110}{kJ}$ por mol de electrones.
\textbf{5.} Dos fotones: unos $176 + 171 = \qty{347}{kJ}$; almacenados,
\qty{110}{kJ}: un \qty{32}{\%}.
\textbf{6.} Cuatro electrones por $\mathrm{O_2}$ (dos moléculas de agua),
que dan dos NADPH: ocho fotones.
\textbf{7.} 12 NADPH: 24 electrones, 48 fotones.
\textbf{8.} $2.303\times 8.314\times 298\times 3/96\,500 = \qty{0.177}{V}$;
$\qty{17.1}{kJ}$ por mol de protones.
\textbf{9.} $4 + 8 = 12$ protones por $\mathrm{O_2}$: 3 \omterm{def:b1:water-small-molecules:atp}{ATP}.
\textbf{10.} $6\times 3 = 18$ \omterm{def:b1:water-small-molecules:atp}{ATP}: exactamente lo que el ciclo necesita con
esta contabilidad; en la práctica el rendimiento se acerca más a 2,6 por
$\mathrm{O_2}$ (\qty{4.7}{H^+} por \omterm{def:b1:water-small-molecules:atp}{ATP} en el \omterm{def:b1:photosynthesis:chloroplast}{cloroplasto}) y el flujo cíclico
en torno al \omterm{def:b1:photosynthesis:zscheme}{fotosistema} I aporta el resto.
\textbf{11.} $4\times 17.1 = \qty{68}{kJ}$ para un \omterm{def:b1:water-small-molecules:atp}{ATP} de \qty{50}{kJ}:
un \qty{73}{\%}.
\textbf{12.} $10^{-5}\,\mathrm{mol/L}\times 10^{-15}\,\mathrm{L}\times
6\times 10^{23} = 6$ protones libres en la luz tilacoidal frente a $10^5$
por segundo a través de las sintasas: el flujo lo transportan protones
unidos a los grupos tamponantes de las \omterm{def:b1:proteins:peptide}{proteínas} y los \omterm{def:b1:lipids:fattyacid}{lípidos} de esa luz,
que los liberan tan deprisa como se van.
\textbf{13.} $18\times 50 + 12\times 220 = 900 + 2640 = \qty{3540}{kJ}$
para \qty{2870}{kJ} almacenados: \qty{670}{kJ} (una quinta parte) disipados
como calor en las reacciones del ciclo, que deben ir cuesta abajo para
funcionar.
\textbf{14.} $48\times 176 = \qty{8450}{kJ}$; $2870/8450 = \qty{34}{\%}$.
\textbf{15.} $0.34\times 0.45\times 0.9 = \qty{13.8}{\%}$.
\textbf{16.} $400\times 0.005\times 36\,000 = \qty{72}{kJ}$; glucosa,
$0.138\times 72\,000/2\,870\,000 = \qty{3.5e-3}{mol} = \qty{0.62}{g}$.
\textbf{17.} Neto, $0.6\times 0.62 = \qty{0.37}{g}$ de glucosa y
$\qty{0.15}{g}$ de carbono ($72/180$).
\textbf{18.} Luz que cae entre las plantas o sobre el suelo; \omterm{def:b1:flowering-plant-organization:organs}{hojas} a la
sombra de otras \omterm{def:b1:flowering-plant-organization:organs}{hojas} trabajando por debajo de su capacidad, y \omterm{def:b1:flowering-plant-organization:organs}{hojas} a pleno
sol saturadas (las \omterm{def:b1:enzymes:enzyme}{enzimas} no pueden usar todos los fotones); respiración de
raíces, \omterm{def:b1:flowering-plant-organization:organs}{tallos} y de noche; fotorrespiración; las semanas anteriores al
cierre de la cubierta y posteriores a su senescencia.
\textbf{19.} 33 oxigenaciones, 16,7 carbonos perdidos, 83,3 ganados netos.
\textbf{20.} Rescate, $33\times 3.5 = 117$ \omterm{def:b1:water-small-molecules:atp}{ATP}; ciclo, 300 \omterm{def:b1:water-small-molecules:atp}{ATP}; en total,
417 \omterm{def:b1:water-small-molecules:atp}{ATP} para 83,3 carbonos: 5,0 \omterm{def:b1:water-small-molecules:atp}{ATP} por carbono neto.
\textbf{21.} $3 + 2 = 5$ \omterm{def:b1:water-small-molecules:atp}{ATP} por carbono, sin perder carbono: igualdad en
\omterm{def:b1:water-small-molecules:atp}{ATP}, pero la planta C3 ha perdido además el \qty{17}{\%} de su capacidad de
carboxilación; a \qty{25}{\celsius}, prácticamente un empate.
\textbf{22.} 50 oxigenaciones, 25 carbonos perdidos, 75 netos; rescate, 175
\omterm{def:b1:water-small-molecules:atp}{ATP}; $475/75 = 6.3$ \omterm{def:b1:water-small-molecules:atp}{ATP} por carbono neto frente a los 5 de la \omterm{def:b1:photosynthesis:c4cam}{planta C4}, y
con la cuarta parte del carbono perdida: la C4 gana con claridad.
\textbf{23.} Triplicar el $\mathrm{CO_2}$ desplaza la competencia en la
\omterm{def:b1:photosynthesis:fixation}{RuBisCO} hacia la carboxilación, con lo que recorta la fotorrespiración y
eleva la fijación neta en las plantas C3; las \omterm{def:b1:photosynthesis:c4cam}{plantas C4} ya saturan la
\omterm{def:b1:photosynthesis:fixation}{RuBisCO} con $\mathrm{CO_2}$ concentrado y no ganan más que un pequeño ahorro
de agua.
\textbf{24.} \qty{1}{g} de carbono son \qty{0.083}{mol}: C3, $250\times 0.083 = \qty{21}{mol}$ de agua, \qty{375}{g}; CAM, \qty{37}{g}.
\textbf{25.} Unos 8 fotones por $\mathrm{CO_2}$ (48 por glucosa); un
\qty{34}{\%} desde los fotones absorbidos hasta la glucosa; alrededor de un
\qty{14}{\%} desde la luz solar incidente hasta el azúcar bruto y un
\qty{8}{\%} hasta el azúcar neto de la \omterm{def:b1:flowering-plant-organization:organs}{hoja}, cifra que las pérdidas de una
planta entera a lo largo de toda una temporada reducen al \qty{1}{\%}.
\end{solution}
